a) Batas
\mathl{\partial M}{} suatu manifold berbatas, menurut
[[Mannigfaltigkeit mit Rand/Rand/Elementare Eigenschaften/Fakt|Fakta (2)]],
\definitionsverweis {tertutup}{}{,}
sedangkan interval terbuka sebagai himpunan bagian dari $\R^2$ tidak tertutup.

b) Batas
\mathl{\partial M}{} suatu manifold berbatas, menurut
[[Mannigfaltigkeit mit Rand/Rand ist Mannigfaltigkeit/Fakt|Fakta]],
merupakan manifold tanpa batas, sedangkan interval tertutup memiliki titik-titik batas.

c) Andaikan bahwa $\R^2$ merupakan manifold dengan batas
\mathl{\R = \R \times \{0\}}{.} Untuk
\mathl{P \in \R}{} terdapat suatu lingkungan terbuka
\mathl{P\in U}{} dan suatu homeomorfisme
\maabbdisp {\alpha} { U } {V
} {}
dengan suatu himpunan terbuka $V \subseteq H$, dengan $H$ menyatakan suatu setengah bidang. Dalam hal ini, kita dapat mengganti $V$ dengan suatu lingkungan setengah bola terbuka yang lebih kecil di sekitar $P$. Menurut
[[Mannigfaltigkeit mit Rand/Rand/Elementare Eigenschaften/Fakt|Fakta (2)]],
pada peta seperti itu titik-titik batas dipetakan ke titik-titik batas, yaitu

\mavergleichskettedisp
{\vergleichskette
{\alpha ( U \cap \R  )
}
{ =} { V \cap \partial H
}
{ } {
}
{ } {
}
{ } {
}
}
{}{}{.}
Dengan demikian, kita juga memperoleh suatu homeomorfisme antara komplemen-komplemennya, yakni antara
\mathkor {} {U \setminus U \cap \R} {dan} {V \setminus V \cap \partial H} {.}
Namun, lingkungan setengah bola di ruas kanan
\definitionsverweis {terhubung lintasan}{}{,}
sedangkan himpunan di ruas kiri merupakan gabungan saling lepas dari irisannya dengan separuh positif dan separuh negatif. Kedua irisan tersebut terbuka dan juga tidak kosong karena $U$ memuat suatu lingkungan bola dari $P$. Oleh karena itu, himpunan di ruas kiri tidak terhubung, sehingga tidak mungkin terdapat suatu homeomorfisme.

d) Sama seperti c).

 [[Kategorie:Latexseite]]
